\documentclass[11pt]{article}

\usepackage[margin=1in]{geometry}
\usepackage{amsmath,amssymb,amsthm,mathtools}
\usepackage{booktabs}
\usepackage{enumitem}
\usepackage{microtype}
\usepackage[colorlinks=true,linkcolor=blue,citecolor=blue,urlcolor=blue]{hyperref}

\theoremstyle{plain}
\newtheorem{theorem}{Theorem}[section]
\newtheorem{proposition}[theorem]{Proposition}
\newtheorem{lemma}[theorem]{Lemma}
\newtheorem{corollary}[theorem]{Corollary}

\theoremstyle{definition}
\newtheorem{definition}[theorem]{Definition}

\theoremstyle{remark}
\newtheorem{remark}[theorem]{Remark}

\newcommand{\GammaSet}{\Gamma_{\mathrm{set}}}
\newcommand{\GammaEff}{\Gamma_{\mathrm{eff}}}
\newcommand{\GammaB}{\Gamma_{\mathcal B}}
\newcommand{\eps}{\varepsilon}

\title{Algorithms and Certificates for Finite Effective Causal Descent:\\
Width-Based Synthesis for Sensor-Limited Mechanical Control}
\author{Anonymous}
\date{}

\begin{document}
\maketitle

\begin{abstract}
Finite decision systems under partial observation face a simple obstruction: a
controller cannot choose different actions in scenarios that its sensors do not
distinguish.  This article develops the finite algorithmic companion to
effective causal descent.  The focus is not a broader philosophical taxonomy,
but a proof-producing synthesis pipeline for finite sensor-limited mechanical
models.  A finite presentation consists of policy variables, finite domains, and
explicit local relations; global policies are exactly the assignments satisfying
all relations.  Effectivity and bounded realisability are represented by finite
implementation catalogues, costs, and budgets.  The main algorithmic layer is a
certificate-producing dynamic program over tree decompositions, with ordinary
and weighted-width bounds, independently checkable message tables, and typed
certificates for policies, refutations, optimum costs, fibre conflicts, sensor
repairs, finite-element rows, and sparsification ladders.  The mechanics layer
turns certified affine response rows into inner and outer finite ECD models,
proving the sound sandwich
\(\Gamma(P^-)\subseteq\Gamma(P^{\mathrm{true}})\subseteq\Gamma(P^+)\).
Width-adaptive sparsification gives monotone certified refinement rather than a
performance claim.  An exact two-spring demonstrator and a finite prototype
artifact exercise the certificate grammar, reject corrupted certificates, and
produce controlled-width benchmark reports.  External solver output is treated
only as auxiliary evidence unless cross-checked by native ECD certificates;
timeouts and missing solvers remain inconclusive.
\end{abstract}

\noindent\textbf{Keywords:} finite effective causal descent; sensor-limited control;
certifying algorithms; tree decompositions; weighted width; finite-element
certification; sparsification; independent verification.

\tableofcontents

\section{Introduction}
\label{sec:introduction}

A mechanically admissible response need not be implementable by a controller
whose observations fail to distinguish the situations requiring different
actions.  A load case, damage state, or inspection record may determine which
actions are safe, while the available sensor readings merge several such cases
into one information state.  Once the information state is fixed, the controller
has only one command available.  The first task is therefore not to solve a
continuous mechanics problem in full generality, but to certify the finite
relation between observations, admissible actions, and implementable policies in
a declared model.

This article develops the finite, proof-producing side of effective causal
descent (ECD).  The foundational monograph supplies the presentation-relative
language: local validity, compatibility of a diagram, effective realisability,
and bounded realisability.  Here we specialize that language to explicit finite
instances and to mechanically generated finite relations.  The guiding rule is:
\emph{topological order is the constraint; feasibility is the priority}.  We
first fix the finite input representation, then the certificate grammar, then
the solver and verifier, then the mechanical relation generator, then the
sparsification and benchmark layers.  Comparative performance claims come last
and require experiments; they are not consequences of the structural theory.

The contribution is organized around certificates.  A successful run should
return a policy that an independent verifier checks.  A negative run should
return a refutation certificate, not merely a solver log.  A mechanical row
should carry the numerical or exact-arithmetic evidence justifying its safe,
unsafe, or unresolved status.  A sparsified instance should come with interval
bounds proving the appropriate inner--outer inclusions.  The computational
statuses are deliberately limited: certified satisfiable, certified
unsatisfiable, certified optimal, or inconclusive.

\section{Finite ECD presentations}
\label{sec:finite-presentations}

\begin{definition}[Finite ECD presentation]
A finite ECD presentation is a tuple
\[
  \mathcal P=(X,(D_x)_{x\in X},(S_i,R_i)_{i\in M}),
\]
where \(X\) is a finite set of policy variables, each \(D_x\) is a finite
nonempty domain, each scope \(S_i\subseteq X\) is finite, and
\[
  R_i\subseteq\prod_{x\in S_i}D_x
\]
is an explicitly represented permitted table.  Its global policy space is
\[
  \Gamma(\mathcal P)=
  \{a\in\prod_{x\in X}D_x: a|_{S_i}\in R_i\text{ for all }i\in M\}.
\]
\end{definition}

In a sensor-limited mechanics model, a variable often has the form \((j,o)\),
where \(j\) is an actuator or decision channel and \(o\) is an information state
produced by the available sensors.  The finite relation tables record which
joint commands are certified safe for each scenario or local mechanical
constraint.

\begin{definition}[Finite implementation catalogue]
For controller \(j\), let \(\mathcal L_j\) be a finite catalogue of
implementations with evaluation map \(\operatorname{eval}_j(\ell,o)\) and cost
\(c_j(\ell)\).  Introduce implementation variables
\(\ell_j\in\mathcal L_j\cup\{\bot\}\), where \(\bot\) is compatible with every
behavioural assignment and has infinite cost.  A finite-cost catalogue policy is
one satisfying
\[
  x_{j,o}=\operatorname{eval}_j(\ell_j,o)
  \qquad(o\in O_j)
\]
with every \(\ell_j\ne\bot\).
\end{definition}

Thus finite set-theoretic existence is separated from existence in a declared
implementation catalogue, and budgeted existence is separated again by a cost
bound.  This is the finite analogue of the nested spaces
\(\GammaSet\supseteq\GammaEff\supseteq\GammaB\), but it is entirely explicit.

\section{Certificate grammar and status discipline}
\label{sec:certificates}

The verifier accepts typed certificates rather than solver narratives.  The
current grammar contains the following core forms.

\begin{center}
\begin{tabular}{@{}p{0.23\linewidth}p{0.41\linewidth}p{0.25\linewidth}@{}}
\toprule
Certificate & Data & Certified claim\\
\midrule
\texttt{POLICY} & total assignment to all variables & global section\\
\texttt{FIBRE-CONFLICT} & one observation fibre and excluded-action witnesses & no one-shot selector on that fibre\\
\texttt{SENSOR-REPAIR} & selected sensors and conflict cores & sensor sufficiency or optimality in a finite master problem\\
\texttt{TREE-DECOMP} & bags, edges, root, factor placement & valid decomposition\\
\texttt{DP-REFUTATION} & decomposition and message tables & no global assignment\\
\texttt{DP-OPTIMUM} & decomposition, message tables, root value & minimum finite cost\\
\texttt{FE-ROW} & response interval and row status & safe, unsafe, or unresolved finite row\\
\texttt{SPARSIFICATION-BATCH} & retained scopes and omitted-term bounds & inner--outer sparse model\\
\bottomrule
\end{tabular}
\end{center}

\begin{remark}[Status discipline]
The implementation reports only four mathematical statuses:
\texttt{CERTIFIED-SAT}, \texttt{CERTIFIED-UNSAT}, \texttt{CERTIFIED-OPT}, and
\texttt{INCONCLUSIVE}.  A timeout, missing external solver, failed certificate
verification, or unresolved numerical row is inconclusive.  It is never a proof
of infeasibility.
\end{remark}

\section{One-shot fibres and causal compilation}
\label{sec:fibres-compilation}

\begin{theorem}[One-shot fibre criterion]
Let \(o:W\to O\) be an observation map, let \(A\) be a finite action set, and
let \(S_w\subseteq A\) be the actions admissible in scenario \(w\).  An
observation-uniform action selector exists if and only if
\[
  A_z:=\bigcap_{w:o(w)=z}S_w\ne\varnothing
\]
for every realised observation value \(z\).  If a fibre fails, a certificate
using at most \(|A|\) scenarios in that fibre excludes every action, and the
bound is tight.
\end{theorem}

\begin{proof}
A selector chooses one action \(a_z\) for every observation value.  The action is
valid on the fibre over \(z\) exactly when \(a_z\in S_w\) for every scenario in
that fibre, equivalently \(a_z\in A_z\).  If \(A_z=\varnothing\), choose for each
action \(a\in A\) one scenario \(w_a\) in the fibre with \(a\notin S_{w_a}\).
The distinct chosen scenarios exclude all actions and have cardinality at most
\(|A|\).  Tightness is witnessed by \(S_{w_j}=A\setminus\{j\}\) for
\(A=\{1,\ldots,k\}\).
\end{proof}

\begin{theorem}[Reachability-correct forbidden-tuple compilation]
Consider a finite explicit history tree whose predictor nodes are partitioned
into information cells.  Introduce one policy variable for each information
cell.  For each unsafe leaf, read the information-cell/action labels along the
root-to-leaf path.  If the path requires the same information cell to take two
different actions, discard it.  Otherwise add the corresponding forbidden tuple.
The resulting finite presentation has global assignments exactly the policies
that avoid every unsafe leaf.
\end{theorem}

\begin{proof}
A path requiring two different actions at one information cell is unrealisable
by any observation-uniform policy, so it contributes no feasible unsafe
execution.  Every remaining unsafe path is reached precisely when the policy
matches all action labels on that path; forbidding the resulting tuple is
therefore exactly the condition that this unsafe leaf is not reached.
\end{proof}

\section{Finite complexity boundary}
\label{sec:complexity}

The explicit-table representation fixes the complexity claims.  Verifying a
proposed assignment is polynomial by scanning factor rows.  Unary constraints
are solved by intersecting allowed values.  Boolean binary constraints are
polynomial via 2-SAT.  Boolean ternary explicit-table CSP is NP-complete by the
standard 3-SAT encoding, and binary disequality over three-element domains is
NP-complete by graph 3-colouring.  Nonexistence is coNP-complete, and counting
global assignments is \(\#\mathrm P\)-complete.

These are not claims about arbitrary succinct generators.  Circuits, automata,
and finite-element preprocessors change the input representation and require
separate analysis.

\section{Width-based certifying dynamic programming}
\label{sec:width-dp}

Let \((T,(B_t)_{t\in T})\) be a tree decomposition of the primal graph.  The
solver orients \(T\) at a root, assigns each factor to a containing bag, and
computes message tables from leaves to root.  A message entry summarizes the
best extension of a boundary assignment through the processed subtree.

For heterogeneous domains, the relevant state-space parameter is the weighted
width
\[
  \lambda(T)=\max_{t\in T}\sum_{x\in B_t}\lceil\log_2|D_x|\rceil.
\]
Every bag table has at most \(2^{\lambda(T)}\) entries.  Thus the explicit-table
recurrences run in \(\operatorname{poly}(|I|,|T|)2^{\lambda(T)}\), up to
indexing and arithmetic factors.

\begin{proposition}[DP certificate soundness]
If an independent verifier recomputes all message tables of a
\texttt{DP-REFUTATION} certificate and obtains an empty root table, then
\(\Gamma(\mathcal P)=\varnothing\).  If it recomputes a \texttt{DP-OPTIMUM}
certificate with root value \(c\), then there is a global assignment of cost
\(c\), and no lower-cost assignment exists for the declared finite instance.
\end{proposition}

\begin{proof}
By induction on the rooted decomposition, each message table is the exact
aggregate of all assignments in the processed subtree agreeing with the boundary
assignment.  Projection combines alternatives by minimization and combines
independent subtree costs by addition.  The root has no boundary, so its table
states exactly whether any global assignment exists and, when costs are present,
the minimum cost.
\end{proof}

\section{Sensor repair}
\label{sec:sensor-repair}

A sensor architecture is part of the presentation.  For a one-shot conflict core
\(C\subseteq W\), let
\[
  E_C=\{s:s\text{ is nonconstant on }C\}
\]
be the candidate sensors that separate the scenarios in the core.

\begin{theorem}[One-shot sensor repair]
A selected sensor family \(T\) repairs all listed one-shot conflict cores if and
only if
\[
  T\cap E_C\ne\varnothing
\]
for every listed core \(C\).  If the list contains all minimal failing fibres,
this hitting-set condition is exact for one-shot repair.
\end{theorem}

The theorem separates two tasks: synthesizing a policy for a fixed observation
architecture, and designing an observation architecture that makes policy
synthesis possible.  The second task contains hitting-set structure even when
the first task is easy.

\section{Certified affine mechanics front end}
\label{sec:mechanics-front-end}

The mechanics front end supplies finite relations; the finite solver certifies
policies or refutations for those relations.  For affine response rows,
\[
  y_i(a)=d_i+\sum_{x\in S_i}g_{ix}a_x,
\]
exact arithmetic or certified intervals classify each row as safe, unsafe, or
unresolved.  More generally, for a symmetric positive-definite discrete solve
\(Ku=f\), residual \(r=f-K\widehat u\), and certified lower spectral bound
\(\underline\lambda\le\lambda_{
\min}(K)\), one has
\[
  \|u-\widehat u\|_2\le \frac{\|r\|_2}{\underline\lambda},
  \qquad
  |c^T(u-\widehat u)|\le
  \|c\|_2\frac{\|r\|_2}{\underline\lambda}.
\]
Rows whose certified response intervals cross a threshold remain unresolved.

\begin{proposition}[Inner--outer finite soundness]
If finite row certificates give
\[
  R_i^-\subseteq R_i^{\mathrm{true}}\subseteq R_i^+
\]
for every factor, then
\[
  \Gamma(P^-)\subseteq\Gamma(P^{\mathrm{true}})\subseteq\Gamma(P^+).
\]
Thus inner satisfiability gives a certified true finite-model policy, and outer
unsatisfiability gives a certified true finite-model refutation.  Inner
unsatisfiability with outer satisfiability is inconclusive.
\end{proposition}

\section{Width-adaptive sparsification}
\label{sec:sparsification}

For affine rows, retain \(J_i\subseteq S_i\) and bound omitted terms by
\[
\underline\eta_i(J)=
\sum_{x\in S_i\setminus J_i}\min_{v\in D_x}g_{ix}v,
\qquad
\overline\eta_i(J)=
\sum_{x\in S_i\setminus J_i}\max_{v\in D_x}g_{ix}v.
\]
The inner sparse relation keeps a retained assignment only if every completion
of omitted variables is certified safe.  The outer sparse relation keeps it if
some completion is not certified unsafe.

\begin{theorem}[Sparsification ladder]
If retained scopes grow from \(J_i\) to \(J_i'\), then
\[
  \Gamma^-_J
  \subseteq
  \Gamma^-_{J'}
  \subseteq
  \Gamma^{\mathrm{true}}
  \subseteq
  \Gamma^+_{J'}
  \subseteq
  \Gamma^+_J.
\]
\end{theorem}

This is a soundness theorem, not a performance theorem.  It says that refinement
is monotone and certifiable; whether it improves runtime on a class of mechanics
instances is an empirical question.

\section{Prototype artifact and reproducibility}
\label{sec:artifact}

The current prototype is located at
\begin{verbatim}
/home/user/finite_ecd_prototype/
\end{verbatim}
It includes explicit-table solving, tree-decomposition DP certificates, exact
affine row certificates, sparsification certificates, controlled-width benchmark
generation, baseline encoders, optional external-solver adapters, external-result
cross-checking, report summaries, and SVG plot generation.  The repository-ready
entry point is
\begin{verbatim}
make reproduce
\end{verbatim}
with \texttt{./reproduce.sh} as an equivalent shell command.

At the present stage, no repository URL or DOI has been supplied.  The paper
should therefore use placeholder wording: an artifact containing the finite ECD
prototype and reproduction scripts will be deposited in a public repository
before publication.  The release manifest should be frozen only after deposit.

\section{Limitations and conclusions}
\label{sec:limitations}

The companion theory is deliberately finite and certificate-driven.  It does not
claim a universal dichotomy for all finite ECD instances, practical superiority
over generic solvers, or continuum-mechanics validity from finite algebraic row
certificates alone.  Its claim is narrower and checkable: once a finite source
representation is declared, the solver and verifier can produce typed evidence
for policies, refutations, optimum costs, sensor repairs, finite mechanical
relations, and sparsification ladders.  The location of failure remains part of
the presentation, and the certificate records the data needed to transport that
claim.

\appendix
\section{Artifact, certificate formats, and reproducibility}
\label{app:artifact}

This appendix describes the finite ECD prototype accompanying the finite
algorithmic and mechanics development. The artifact is a correctness and
certificate harness for finite explicit-table instances and exact finite affine
algebraic examples. It is not a continuum-mechanics validator and it is not, by
itself, a performance comparison with SAT, CP-SAT, MILP, or other external
solvers.

\subsection{Finite instance format}

A finite explicit-table ECD instance is encoded as a JSON object with the fields
\texttt{variables}, \texttt{domains}, and \texttt{factors}. A typical Boolean
example is
\begin{verbatim}
{
  "variables": ["x0", "x1"],
  "domains": {"x0": [0, 1], "x1": [0, 1]},
  "value_costs": {"x0": {"0": 0, "1": 1}},
  "factors": [
    {"name": "neq", "scope": ["x0", "x1"],
     "relation": [[0, 1], [1, 0]]}
  ]
}
\end{verbatim}
The field \texttt{variables} fixes the ordered policy variables; \texttt{domains}
gives each finite domain; \texttt{factors} gives explicit permitted tuples on
ordered scopes; and \texttt{value\_costs}, when present, supplies nonnegative
value costs for optimum certificates. The complexity and certificate claims made
for this artifact are explicit-table claims unless a different generator is
specified.

\subsection{Native certificate schemas}

The verifier recognises several certificate types. A \texttt{POLICY} certificate
assigns one domain value to each variable and is checked by scanning all factor
relations. A \texttt{FIBRE-CONFLICT} certificate names an observation fibre and,
for each action, a scenario in that fibre excluding the action. A
\texttt{SENSOR-REPAIR} certificate records a selected sensor set for the small
one-shot repair examples; the current optimality certificates are exhaustive
certificates for tiny instances.

A \texttt{TREE-DECOMP} certificate contains bags, tree edges, a root bag, and a
factor-to-bag map:
\begin{verbatim}
{
  "type": "TREE-DECOMP",
  "root": "b0",
  "bags": {"b0": ["x0", "x1"]},
  "edges": [],
  "factor_to_bag": {"neq": "b0"}
}
\end{verbatim}
The verifier checks factor containment, tree connectivity, and the
running-intersection property. A \texttt{DP-OPTIMUM} or \texttt{DP-REFUTATION}
certificate contains a tree decomposition, dynamic-programming message tables,
and a root optimum or refutation. The verifier recomputes all messages; solver
logs are not treated as certificates.

For exact affine mechanics-style examples, an \texttt{FE-ROW} certificate records
the row name, action assignment, exact response interval, safety interval, and
classification as \texttt{SAFE} or \texttt{UNSAFE}. An \texttt{FE-ROW-BATCH}
collects such rows and includes the derived finite ECD relation-table instance.
A \texttt{SPARSIFICATION-BATCH} records retained scopes, exact omitted-term
intervals, retained-response intervals, and the corresponding inner and outer
sparse finite instances. A sparsification-ladder check compares two such batches
for retained scopes $J\subseteq J'$: coarse inner cylinders must refine into fine
inner rows, and fine outer rows must project into coarse outer rows.

\subsection{Baseline encodings and external solvers}

The script \texttt{baseline\_encode.py} emits DIMACS CNF, LP-format MILP, and a
minimal CP-SAT-style forbidden-tuple JSON encoding for Boolean explicit-table
instances. These encodings are supplied to make future comparisons use identical
finite instances. They are not solver results.

External solver output is handled conservatively. The optional solver adapters
record availability, timeouts, and raw status. The cross-checker accepts an
external SAT or UNSAT result only when an independent native ECD certificate is
also verified: SAT requires a policy certificate, and UNSAT requires a verified
\texttt{DP-REFUTATION}. Missing solvers, timeouts, and unrecognised output are
reported as \texttt{INCONCLUSIVE}.

\subsection{Reproduction commands}

From the artifact directory, the repository-ready entry point is
\begin{verbatim}
make reproduce
\end{verbatim}
or equivalently
\begin{verbatim}
./reproduce.sh
\end{verbatim}
This runs the smoke checks, controlled-width benchmark sweep, and SVG plot
generation. The principal outputs are \texttt{repro\_report.json},
\texttt{reports/benchmark\_sweep.json},
\texttt{reports/certificate\_size\_vs\_width.svg}, and
\texttt{reports/verify\_time\_vs\_certificate\_size.svg}.

\subsection{Claim boundary}

The artifact certifies finite explicit-table examples and exact finite affine
algebraic models. It does not certify continuum discretisation error, physical
modelling adequacy, or practical superiority over external solvers. External
baseline experiments should be reported only after solver versions, hardware,
time limits, memory limits, and encodings have been fixed; timeouts remain
\texttt{INCONCLUSIVE}.


\end{document}
